% New finite-cohort negotiation analysis. Canonical inputs unchanged.
% See negotiation/results.json and protocol.md. Semantic counts are one
% model-assisted manual pass pending independent adjudication.

\subsection{A complete negotiation population}
\label{sec:agenda-negotiation-method}

We reconstruct each negotiation from its initial proposal through its terminal
acceptance or rejection. A counteroffer creates a new response opportunity
within that thread. The actor issuing the current offer owns the offered
bundle; the original thread initiator need not do so. We cross-check that
orientation against the recipient's delivered \texttt{if\_you\_accept} fields,
then calculate immediate cash, mortgage-transfer interest, and newly completed
or relinquished color groups for both parties. These calculations characterize
the available exchange rather than assigning it a strategic utility. They
exclude railroads and utilities from the color-completion measure and retain
mortgaged deeds when testing ownership completion.

The canonical events contain 526 complete threads and 920 response
opportunities: 113 acceptances, 413 rejections, and 394 counters. Acceptance is
offered by the engine in 918 opportunities. The two exceptions would leave one
party with negative cash after mandatory charges, so they are excluded from
comparisons requiring a genuine choice to accept. The accounting calculation
agrees with the delivered acceptance menu in all 920 cases. Reconstructed cash,
ownership, mortgage and development state agrees with all 6,989 initial decision
prompts, and all 113 predicted accepted-transfer states agree with the next
recorded state. A separate checker recomputes group completion directly from
the delivered bundles and rejects eight deliberately corrupted outputs.

The protocol fixes both the population and its units. A settlement fraction
uses threads as its denominator; a response fraction uses offers received.
No negotiation spans a new global game turn in this implementation: the engine
resolves the exchange before normal play resumes. The longest thread contains
eleven counters, and the median contains one response. Of 526 threads, 414
follow an earlier thread between the same pair in the same game. Repeated
offers, shared positions and the five seed blocks therefore preclude treating
920 responses as independent trials. We provide all eighty game/configuration
cells and all 120 game/pair cells, including zeros.

\subsection{Completing one's own set is insufficient for acceptance}
\label{sec:agenda-completion-response}

Table~\ref{tab:agenda-completion} gives the complete response distribution by
the immediate ownership consequences of accepting. Among 221 legal offers that
would newly complete a recipient color group, 34 are accepted, 113 are countered
and 74 are rejected. Even when both parties would complete a group, only ten
of 81 offers are immediately accepted; thirty-five are countered and thirty-six
rejected. These observations answer a bounded version of the strategic-reasoning
question: obtaining a new monopoly is not a sufficient condition for agreement
in the recorded configured systems. They do not show that each refusal maximizes
winning probability, or that all recipients correctly assess the alternatives.

\begin{table}[t]
\centering\small
\begin{tabular}{lrrrr}
\toprule
New color completion if accepted & Offers & Accept & Counter & Reject\\
\midrule
Neither party & 408 & 60 & 169 & 179\\
Current author only & 289 & 19 & 111 & 159\\
Recipient only & 140 & 24 & 78 & 38\\
Both parties & 81 & 10 & 35 & 36\\
\midrule
All acceptance-legal offers & 918 & 113 & 393 & 412\\
\bottomrule
\end{tabular}
\caption{Response to every acceptance-legal offer. ``New completion'' is an
ownership fact, not a declaration of buildability or gain in utility. Two of
the 920 received offers lack acceptance in the legal menu and are excluded:
one is rejected and one countered. The same thread and player can contribute
several rows.}
\label{tab:agenda-completion}
\end{table}

Own completion can also require sacrificing an existing set. Six of the 221
offers do so; one of the 74 direct rejections is in this category. Removing
that rejection still leaves 73 refusals that newly complete a recipient group
without breaking another complete recipient group. Other costs remain:
cash, mortgage charges, timing, development expense and the rival's new
position. Accordingly, the test refutes a sufficient-condition interpretation
of set completion, rather than demonstrating that any particular trade was
beneficial overall.

To examine the stated considerations, we read the complete applied public
message and private note for all 74 rejections in this closed population.
The nonexclusive manual codes record rival completion, development, financing,
survival or estate denial in 63; own price or liquidity in 42; portfolio
retention in 21; and explicit timing or a reconsideration condition in 13.
All 74 have public text and 66 have a private note. These counts describe
expressed considerations in those two channels, not verified calculations,
latent motives or their causal influence. They are a single model-assisted
coding pass with exact source spans; independent adjudication is required
before treating the semantic counts as validated labels. Repeated negotiations
concentrate the evidence: G1 contributes 22 of the 74 and G15 contributes twelve.

The exceptions matter. At G15 d106 and G17 d186, the purported justification
adopts the seller's role although the recorded actor is the buyer. These two
records are flagged rather than interpreted as coherent reservation prices.
Two other declines contain no applicable reason in the two coded fields;
this does not establish the absence of a rationale in another channel.
No response in the 920-opportunity population is an engine default, but 168
use a retry. Opus contributes 162 action-only responses among its 189 response
opportunities, so differences in explanation availability remain configuration
conditions. A valid received action also need not match accompanying content.

\subsection{Negotiation activity and settlement are different outcomes}
\label{sec:agenda-negotiation-outcomes}

Only 32 settlements occur without a counteroffer. Eighty-one follow at least
one counter, and in 44 settlements the final author differs from the initial
proposer. Of the 388 threads without counters, 32 settle; of the 138 with
counters, 81 settle. This association does not establish that countering
improves a deal: participants select which conversations to continue, and each
counter changes the terms.

The initiating and responding roles also produce different descriptions.
Sol initiates 132 threads, of which nineteen settle; Gemini initiates eighteen,
of which eleven settle; Opus initiates 240, of which forty-four settle; and
Grok initiates 136, of which thirty-nine settle. As responders, their acceptance
counts are respectively thirty, twenty-six, twenty-nine and twenty-eight out
of 235, 242, 189 and 254 received offers. These are exposure- and role-specific
frequencies, not negotiation-skill rankings. Gemini's low initiation frequency
and high conditional settlement fraction could reflect selective proposals,
favorable inventory, the opponents' demands or several of these together.

The complete ledger supplies an answer about how these participants negotiate,
but does not establish payoff quality, agreement bias or persuasion. A price
above the bank list price is not necessarily overpayment, and an accepted
exchange completing both sides is not necessarily fair. A future utility test
must specify the outside option and continuation distribution, including effects
on third parties. The contextual transcripts complement these population counts
by showing what was said and received; their selected nature remains explicit.

\subsection{Collective outcomes through distinct transfer mechanisms}
\label{sec:agenda-collective-flows}

We construct separate directed networks for trade cash, traded deeds, paid
rent, estate cash and inherited deeds. Currency and deed counts are never
added into one edge weight. Actual settlement transfers comprise 107 directed
cash flows totaling \$28,425 and 207 deed transfers. There are 905 paid-rent
events totaling \$121,154. Estate cash totals \$11,830 across 59 nonzero
transfers, including \$16 paid to the bank; 397 deeds pass to player-creditors.
The estate figures include repeated inheritance of the same property, and
unpaid terminal rent demands are excluded from paid rent.

The eventual winner receives \$78,737 of paid rent and 372 of the 397 inherited
deeds. Terminal estates mechanically concentrate assets in the remaining
player, so we report a separate nonterminal partition: winners receive 202 of
227 inherited deeds and \$7,148 of \$7,843 estate cash before the final
elimination. The remaining 170 inherited deeds are terminal transfers. These
data describe realized concentration through creditor inheritance. They do not
estimate how much inheritance causes victory, because creditor identity,
development, cash and landings evolve jointly.

The ledger also retains deed-share concentration at all 6,989 decision requests
and every bilateral flow, permitting game-specific trajectories without mixing
cash and ownership. A dense trading pair or concentrated flow is not evidence
of collusion. The existing coordination transcript establishes an invitation
and observed responses, while a claim about coalition incidence or stability
requires a defined episode, acknowledgment and follow-through codebook. The
new population analysis therefore answers part of the collective-dynamics
question through measured transfer mechanisms; it does not manufacture a
coalition from ordinary interdependence.
